\documentclass[UTF8,a4paper,zihao=-4]{ctexart}

\usepackage[left=2.5cm,right=2.5cm,top=2.6cm,bottom=2.6cm]{geometry}
\usepackage{amsmath,amssymb}
\usepackage{booktabs}
\usepackage{array}
\usepackage{multirow}
\usepackage{float}
\usepackage{siunitx}
\usepackage{enumitem}
\usepackage{setspace}
\usepackage{fontspec}
\usepackage{graphicx}

\setmainfont{Times New Roman}
\setCJKmainfont{SimSun}
\sisetup{detect-all}

\setstretch{1.5}
\setlength{\parindent}{2em}
\pagestyle{plain}

\AtBeginDocument{\songti}

\begin{document}

% ===== 封面 =====
\begin{titlepage}
    \centering
    \includegraphics[width=0.5\textwidth]{C:/Users/Cgy123456/Desktop/学长的焚决/大一下/大物实验/华南理工大学水印.jpg}
    \vspace{1.5cm}
    
    {\zihao{2}\bfseries 实验报告\par}
    \vspace{1cm}
    {\zihao{-2}\bfseries 实验5.20\quad 弗兰克-赫兹实验\par}
    \vspace{2cm}

    \renewcommand{\arraystretch}{1.8}
    \begin{tabular}{p{4cm}p{8cm}}
        课程名称： & 大学物理实验 \\
        姓名： & 【待补充】 \\
        学号： & 【待补充】 \\
        专业班级： & 【待补充】 \\
        开课学期： & 【待补充】 \\
        实验日期： & 【待补充】 \\
    \end{tabular}

    \vfill
    {\zihao{4} 物理与光电学院公共教学实验中心\par}
    \vspace{0.5cm}
    {\zihao{4} 【待补充】年【待补充】月\par}
\end{titlepage}

\newpage

\section*{实验名称}

实验5.20\quad 弗兰克-赫兹实验

\section*{引言}

原子能级的存在最早是从光谱学的研究中推断出来的。1914年，弗兰克和赫兹采用慢电子和稀薄气体原子碰撞的方法，测得原子吸收或发射的能量是不连续的，从而证明了原子能级的存在。弗兰克和赫兹也由于这一杰出的贡献共同获得了1925年的诺贝尔物理学奖。本实验通过测量氩原子的第一激发电位，验证原子能级的量子化特性。

% 以下正文：宋体，小四号，1.5倍行距

\section*{一、实验目的}

\begin{enumerate}[label=\arabic*.]
    \item 学习测量原子的第一激发电位的方法；
    \item 通过实验证实原子能级的存在；
    \item 研究影响充气电子管阳极电流的因素，分析其机理。
\end{enumerate}

\section*{二、实验仪器}

弗兰克-赫兹实验仪（充气四极电子管）。

\section*{三、实验原理}
% 位置不够，可自行加页

\subsection*{1. 玻尔原子模型}

玻尔关于原子理论的两个基本假设：

（1）定态假设：原子只能较长久地停留在一些定态，原子在定态时不发射也不吸收能量，各定态的能量是不连续的值 \(E_1\)、\(E_2\)、\(E_3\)、…。

（2）跃迁假设：原子从一个定态 \(E_2\) 跃迁到另一个定态 \(E_1\) 时，要辐射出一个光子，其频率是一定的，满足：
\[
h\nu = E_2 - E_1.
\]

\subsection*{2. 第一激发电位}

当电子与原子发生非弹性碰撞时，电子将部分动能传递给原子，使原子从基态跃迁到激发态。设原子从基态 \(E_1\) 跃迁到第一激发态 \(E_2\)，则电子需要提供的能量为
\[
eU_0 = E_2 - E_1,
\]
其中 \(U_0\) 即为原子的第一激发电位。

\subsection*{3. 弗兰克-赫兹实验原理}

在弗兰克-赫兹管中，电子从阴极K发射，经过加速电压 \(U_{G_2K}\) 加速后与氩原子发生碰撞。当电子能量小于第一激发能时，发生弹性碰撞；当电子能量达到或超过第一激发能时，发生非弹性碰撞，电子将能量传递给原子使其激发。

随着加速电压 \(U_{G_2K}\) 的增加，阳极电流 \(I_A\) 呈现周期性变化。相邻两个峰电流（或谷电流）对应的电压差即为原子的第一激发电位 \(U_0\)。

\section*{四、实验内容与步骤}
% 位置不够，可自行加页

\begin{enumerate}[label=\arabic*.]
    \item 按图连接好实验电路，接通电源，开机，预热5分钟。【电子逸出过程与热能有关，需要一定的稳定时间，故需预热！】
    \item 手动、自动测量。
    \item 设置参数：灯丝电压 \(2.2\ \text{V}\)，第一栅极电压 \(1.5\ \text{V}\)，拒斥电压 \(9.0\ \text{V}\)。
    \item "启动/增加 \(U_{G_2K}\) 电压"；
    \item 查阅数据，记录谷底或峰尖对应 \(U_{G_2K}\) 电压，并计算第一激发电位 \(U_0\)。
    \item 按下"自动/手动"键，将数据清零。
\end{enumerate}

\section*{五、数据处理}
% 位置不够，可自行加页

\subsection*{1. 原始数据记录}

\begin{table}[H]
    \centering
    \caption{自动模式测量 \(I_A-U_{G_2K}\) 曲线峰、谷电压记录表}
    \renewcommand{\arraystretch}{1.35}
    \begin{tabular}{ccccccc}
        \toprule
        序号 \(i\) & 1 & 2 & 3 & 4 & 5 & 6 \\
        \midrule
        峰 \(U_{G_2K}\) / V & 22.0 & 32.0 & 43.3 & 55.0 & 67.2 & 79.6 \\
        谷 \(U_{G_2K}\) / V & 26.1 & 37.1 & 48.4 & 60.0 & 72.2 & \\
        \bottomrule
    \end{tabular}
\end{table}

\subsection*{2. 逐差法计算第一激发电位}

由附表5.20.1中峰、谷电压值，用逐差法计算第一激发电势，求其平均值 \(\overline{U_0}\)，并与氩原子第一激发电势公认值 \(U_0 = 11.55\ \text{V}\) 比较，求出相对误差。

\subsubsection*{（1）利用峰电压数据计算}

\[
\begin{aligned}
\Delta U_1 &= U_4 - U_1 = 55.0 - 22.0 = 33.0\ \text{V}, \\
\Delta U_2 &= U_5 - U_2 = 67.2 - 32.0 = 35.2\ \text{V}, \\
\Delta U_3 &= U_6 - U_3 = 79.6 - 43.3 = 36.3\ \text{V}.
\end{aligned}
\]

\[
\overline{\Delta U} = \frac{33.0 + 35.2 + 36.3}{3} = 34.83\ \text{V}.
\]

\[
\overline{U_0} = \frac{\overline{\Delta U}}{3} = \frac{34.83}{3} = 11.61\ \text{V}.
\]

\subsubsection*{（2）利用谷电压数据计算}

\[
\begin{aligned}
\Delta U_1 &= U_4 - U_1 = 60.0 - 26.1 = 33.9\ \text{V}, \\
\Delta U_2 &= U_5 - U_2 = 72.2 - 37.1 = 35.1\ \text{V}.
\end{aligned}
\]

\[
\overline{\Delta U} = \frac{33.9 + 35.1}{2} = 34.50\ \text{V}.
\]

\[
\overline{U_0} = \frac{\overline{\Delta U}}{3} = \frac{34.50}{3} = 11.50\ \text{V}.
\]

\subsubsection*{（3）综合结果}

取峰、谷数据的平均值作为最终结果：
\[
\overline{U_0} = \frac{11.61 + 11.50}{2} = 11.56\ \text{V}.
\]

\subsection*{3. 相对误差计算}

原子第一激发电位公认值 \(U_{\text{理论}} = 11.55\ \text{V}\)，则相对误差为
\[
E_r = \frac{|\overline{U_0} - U_{\text{理论}}|}{U_{\text{理论}}} \times 100\%
= \frac{|11.56 - 11.55|}{11.55} \times 100\%
= 0.09\%.
\]

\section*{六、实验结果与分析}

\subsection*{1. 实验结果}

原子第一激发电位测量值为
\[
\boxed{\overline{U_0} = 11.56\ \text{V}}.
\]

与公认值 \(11.55\ \text{V}\) 比较，相对误差为 \(0.09\%\)。

\subsection*{2. 误差分析}

产生误差的可能原因包括：
\begin{enumerate}[label=\arabic*.]
    \item 仪器读数误差，峰谷位置判断存在主观偏差；
    \item 灯丝电压、第一栅极电压、拒斥电压等参数设置不够精确；
    \item 电子管预热时间不足，温度未完全稳定；
    \item 空间电荷效应和接触电势差的影响；
    \item 氩气纯度及气压波动对测量结果的影响。
\end{enumerate}

\subsection*{3. 实验结论}

本实验通过弗兰克-赫兹实验，测量了氩原子的第一激发电位。实验结果表明，阳极电流随加速电压呈现周期性变化，相邻峰（或谷）之间的电压差基本恒定，验证了原子能级的量子化特性。测得氩原子第一激发电位为 \(11.56\ \text{V}\)，与公认值 \(11.55\ \text{V}\) 非常接近，相对误差仅为 \(0.09\%\)，实验结果良好。

% ===== 封底：原始数据记录 =====
\newpage
\begin{center}
    {\zihao{3}\bfseries 原始数据记录\par}
\end{center}

\vspace{0.5cm}

\begin{table}[H]
    \centering
    \caption{自动模式测量 \(I_A-U_{G_2K}\) 曲线峰、谷电压记录表}
    \renewcommand{\arraystretch}{1.35}
    \begin{tabular}{ccccccc}
        \toprule
        序号 \(i\) & 1 & 2 & 3 & 4 & 5 & 6 \\
        \midrule
        峰 \(U_{G_2K}\) / V & 22.0 & 32.0 & 43.3 & 55.0 & 67.2 & 79.6 \\
        谷 \(U_{G_2K}\) / V & 26.1 & 37.1 & 48.4 & 60.0 & 72.2 & \\
        \bottomrule
    \end{tabular}
\end{table}

\vspace{2cm}

\begin{center}
    \renewcommand{\arraystretch}{2}
    \begin{tabular}{p{6cm}p{6cm}}
        实验者签名：\underline{\hspace{4cm}} & 日期：\underline{\hspace{4cm}} \\
    \end{tabular}
\end{center}

\end{document}
